\documentclass{amsart}
% For dealing with references we use the comment environment
\usepackage{verbatim}
\newenvironment{reference}{\comment}{\endcomment}
%\newenvironment{reference}{}{}
\newenvironment{slogan}{\comment}{\endcomment}
\newenvironment{history}{\comment}{\endcomment}

% For commutative diagrams we use Xy-pic
\usepackage[all]{xy}

% We use 2cell for 2-commutative diagrams.
\xyoption{2cell}
\UseAllTwocells

% We use multicol for the list of chapters between chapters
\usepackage{multicol}

% This is generally recommended for better output
\usepackage{lmodern}
\usepackage[T1]{fontenc}

% For cross-file-references

% Package for hypertext links:
\usepackage{hyperref}

% For any local file, say "hello.tex" you want to link to please

% Theorem environments.
%
\theoremstyle{plain}
\newtheorem{theorem}[subsection]{Theorem}
\newtheorem{proposition}[subsection]{Proposition}
\newtheorem{lemma}[subsection]{Lemma}

\theoremstyle{definition}
\newtheorem{definition}[subsection]{Definition}
\newtheorem{example}[subsection]{Example}
\newtheorem{exercise}[subsection]{Exercise}
\newtheorem{situation}[subsection]{Situation}

\theoremstyle{remark}
\newtheorem{remark}[subsection]{Remark}
\newtheorem{remarks}[subsection]{Remarks}

\numberwithin{equation}{subsection}

% Macros
%
\def\lim{\mathop{\mathrm{lim}}\nolimits}
\def\colim{\mathop{\mathrm{colim}}\nolimits}
\def\Spec{\mathop{\mathrm{Spec}}}
\def\Hom{\mathop{\mathrm{Hom}}\nolimits}
\def\Ext{\mathop{\mathrm{Ext}}\nolimits}
\def\SheafHom{\mathop{\mathcal{H}\!\mathit{om}}\nolimits}
\def\SheafExt{\mathop{\mathcal{E}\!\mathit{xt}}\nolimits}
\def\Sch{\mathit{Sch}}
\def\Mor{\mathop{\mathrm{Mor}}\nolimits}
\def\Ob{\mathop{\mathrm{Ob}}\nolimits}
\def\Sh{\mathop{\mathit{Sh}}\nolimits}
\def\NL{\mathop{N\!L}\nolimits}
\def\CH{\mathop{\mathrm{CH}}\nolimits}
\def\proetale{{pro\text{-}\acute{e}tale}}
\def\etale{{\acute{e}tale}}
\def\QCoh{\mathit{QCoh}}
\def\Ker{\mathop{\mathrm{Ker}}}
\def\Im{\mathop{\mathrm{Im}}}
\def\Coker{\mathop{\mathrm{Coker}}}
\def\Coim{\mathop{\mathrm{Coim}}}

% Boxtimes
%
\DeclareMathSymbol{\boxtimes}{\mathbin}{AMSa}{"02}

%
% Macros for moduli stacks/spaces
%
\def\QCohstack{\mathcal{QC}\!\mathit{oh}}
\def\Cohstack{\mathcal{C}\!\mathit{oh}}
\def\Spacesstack{\mathcal{S}\!\mathit{paces}}
\def\Quotfunctor{\mathrm{Quot}}
\def\Hilbfunctor{\mathrm{Hilb}}
\def\Curvesstack{\mathcal{C}\!\mathit{urves}}
\def\Polarizedstack{\mathcal{P}\!\mathit{olarized}}
\def\Complexesstack{\mathcal{C}\!\mathit{omplexes}}
% \Pic is the operator that assigns to X its picard group, usage \Pic(X)
% \Picardstack_{X/B} denotes the Picard stack of X over B
% \Picardfunctor_{X/B} denotes the Picard functor of X over B
\def\Pic{\mathop{\mathrm{Pic}}\nolimits}
\def\Picardstack{\mathcal{P}\!\mathit{ic}}
\def\Picardfunctor{\mathrm{Pic}}
\def\Deformationcategory{\mathcal{D}\!\mathit{ef}}

\setlength{\emergencystretch}{3em}
\begin{document}
\title[Stacks Project, Tag 0G1F]{328 --- Characteristic-p coordinate modification decreases the nonfinite invariant of an etale algebra lift}
\maketitle
\noindent Copyright (C) 2005--2025 Johan de Jong. Stacks Project authors. Source: \href{https://stacks.math.columbia.edu/tag/0G1F}{Tag 0G1F}.

\noindent Editorial difficulty note (not author text). Difficulty H2: The construction must make an additional generator integral while preserving etaleness and the prescribed quotient map. It chooses a polynomial relation and adds carefully chosen p-power terms to the coordinates, controlling integrality, Kähler differentials and the kernel ideal simultaneously to force strict descent of the invariant..

\noindent Editorial provenance note (not author text). History: excerpt from revision \texttt{a0444\allowbreak 6e57e\allowbreak c1fbc\allowbreak 252a8\allowbreak 71afc\allowbreak ec775\allowbreak 2fb28\allowbreak 07b14}, pione.tex, lines 7473--7565. Reference presentation was adapted; original-author context and core support blocks were retained or added. The mathematical wording is unchanged. Licensed under GNU FDL 1.2 or later; no invariant sections or cover texts. See ../licenses/COPYING. Context blocks reproduce author definitions and supporting results; they are not additional counted problems.

\section{Tricks in positive characteristic}
\label{section-achinger}

\noindent
In Piotr Achinger's paper \href{https://stacks.math.columbia.edu/bibliography}{Bibliography: Achinger} it is shown that an
affine scheme in positive characteristic is always a $K(\pi, 1)$.
In this section we explain the more elementary parts of \href{https://stacks.math.columbia.edu/bibliography}{Bibliography: Achinger}.
Namely, we show that for a field $k$ of positive characteristic
an affine scheme \'etale over $\mathbf{A}^n_k$ is actually finite \'etale
over $\mathbf{A}^n_k$ (by a different morphism). We also show
that a closed immersion of connected affine schemes
in positive characteristic induces an injective map on
\'etale fundamental groups.

\medskip\noindent
Let $k$ be a field of characteristic $p > 0$.
Let
$$
k[x_1, \ldots, x_n] \longrightarrow A
$$
be a surjection of finite type $k$-algebras whose source is the
polynomial algebra on $x_1, \ldots, x_n$. Denote
$I \subset k[x_1, \ldots, x_n]$ the kernel so that we have
$A = k[x_1, \ldots, x_n]/I$. We do not assume $A$ is nonzero
(in other words, we allow the case where $A$ is the zero ring
and $I = k[x_1, \ldots, x_n]$). Finally, we assume given a
finite \'etale ring map $\pi : A \to B$.

\medskip\noindent
Suppose given $k, n, k[x_1, \ldots, x_n] \to A, I, \pi : A \to B$.
Let $C$ be a $k$-algebra. Consider
commutative diagrams
$$
\xymatrix{
& B \\
C \ar[r] & C/\varphi(I)C \ar[u]^\tau \\
k[x_1, \ldots, x_n] \ar[u]^\varphi \ar[r] &
A \ar[u] \ar@/_3em/[uu]_\pi
}
$$
where $\varphi$ is an \'etale $k$-algebra map and $\tau$ is a
surjective $k$-algebra map.
Let $C, \varphi, \tau$ be given. For any $r \geq 0$ and
$y_1, \ldots, y_r \in C$ which generate $C$ as an algebra over
$\Im(\varphi)$ let $s = s(r, y_1, \ldots, y_r) \in \{0, \ldots, r\}$
be the maximal element such that $y_i$ is integral over $\Im(\varphi)$
for $1 \leq i \leq s$. We define $NF(C, \varphi, \tau)$
to be the minimum value of
$r - s = r - s(r, y_1, \ldots, y_r)$ for all choices of
$r$ and $y_1, \ldots, y_r$ as above.
Observe that $NF(C, \varphi, \tau)$ is $0$ if and only if $\varphi$ is finite.



\begin{lemma}
\label{lemma-lower-invariant}
In the situation above, if $NF(C, \varphi, \tau) > 0$, then there exist
an \'etale $k$-algebra map $\varphi'$ and a surjective $k$-algebra map
$\tau'$ fitting into the commutative diagram
$$
\xymatrix{
& B \\
C \ar[r] & C/\varphi'(I)C \ar[u]_{\tau'} \\
k[x_1, \ldots, x_n] \ar[u]^{\varphi'} \ar[r] &
A \ar[u] \ar@/_3em/[uu]_\pi
}
$$
with $NF(C, \varphi', \tau') < NF(C, \varphi, \tau)$.
\end{lemma}

\begin{proof}
Choose $r \geq 0$ and $y_1, \ldots, y_r \in C$ which generate
$C$ over $\Im(\varphi)$ and let $0 \leq s \leq r$ be such that
$y_1, \ldots, y_s$ are integral over $\Im(\varphi)$
such that $r - s = NF(C, \varphi, \tau) > 0$.
Since $B$ is finite over $A$, the image of $y_{s + 1}$ in $B$
satisfies a monic polynomial over $A$. Hence we can find
$d \geq 1$ and $f_1, \ldots, f_d \in k[x_1, \ldots, x_n]$
such that
$$
z = y_{s + 1}^d + \varphi(f_1) y_{s + 1}^{d - 1} + \ldots + \varphi(f_d) \in
J = \Ker(C \to C/\varphi(I)C \xrightarrow{\tau} B)
$$
Since $\varphi : k[x_1, \ldots, x_n] \to C$ is \'etale, we can find a
nonzero and nonconstant polynomial $g \in k[T_1, \ldots, T_{n + 1}]$
such that
$$
g(\varphi(x_1), \ldots, \varphi(x_n), z) = 0
\quad\text{in}\quad
C
$$
To see this you can use for example that
$C \otimes_{\varphi, k[x_1, \ldots, x_n]} k(x_1, \ldots, x_n)$
is a finite product of finite separable field extensions
of $k(x_1, \ldots, x_n)$
(see Algebra, Lemmas \href{https://stacks.math.columbia.edu/tag/00U3}{lemma etale over field (Tag 00U3)})
and hence $z$ satisfies a monic
polynomial over $k(x_1, \ldots, x_n)$. Clearing denominators
we obtain $g$.

\medskip\noindent
The existence of $g$ and
Algebra, Lemma \href{https://stacks.math.columbia.edu/tag/051N}{lemma helper polynomial (Tag 051N)}
produce integers $e_1, e_2, \ldots, e_n \geq 1$ such that
$z$ is integral over the subring $C'$ of $C$ generated by
$t_1 = \varphi(x_1) + z^{pe_1}, \ldots, t_n = \varphi(x_n) + z^{pe_n}$.
Of course, the elements $\varphi(x_1), \ldots, \varphi(x_n)$
are also integral over $C'$ as are the elements
$y_1, \ldots, y_s$. Finally, by our choice of $z$ the element
$y_{s + 1}$ is integral over $C'$ too.

\medskip\noindent
Consider the ring map
$$
\varphi' : k[x_1, \ldots, x_n] \longrightarrow C, \quad
x_i \longmapsto t_i
$$
with image $C'$. Since
$\text{d}(\varphi(x_i)) = \text{d}(t_i) = \text{d}(\varphi'(x_i))$
in $\Omega_{C/k}$ (and this is where we use the characteristic of $k$
is $p > 0$) we conclude that $\varphi'$ is \'etale because $\varphi$ is
\'etale, see
Algebra, Lemma \href{https://stacks.math.columbia.edu/tag/0G1C}{lemma characterize etale over polynomial ring (Tag 0G1C)}.
Observe that $\varphi'(x_i) - \varphi(x_i) = t_i - \varphi(x_i) = z^{pe_i}$
is in the kernel $J$ of the map $C \to C/\varphi(I)C \to B$ by our
choice of $z$ as an element of $J$.
Hence for $f \in I$ the element
$$
\varphi'(f) =
f(t_1, \ldots, t_n) =
f(\varphi(x_1) + z^{pe_1}, \ldots, \varphi(x_n) + z^{pe_n}) =
\varphi(f) + \text{element of }(z)
$$
is in $J$ as well. In other words, $\varphi'(I)C \subset J$ and
we obtain a surjection
$$
\tau' : C/\varphi'(I)C \longrightarrow C/J \cong B
$$
of algebras \'etale over $A$. Finally, the algebra
$C$ is generated by the elements
$\varphi(x_1), \ldots, \varphi(x_n), y_1, \ldots, y_r$
over $C' = \Im(\varphi')$ with
$\varphi(x_1), \ldots, \varphi(x_n), y_1, \ldots, y_{s + 1}$
integral over $C' = \Im(\varphi')$. Hence
$NF(C, \varphi', \tau') < r - s = NF(C, \varphi, \tau)$.
This finishes the proof.
\end{proof}
\end{document}
